\scanpage{60}
\mspage{51}
% archive sheet number: 60
\uline{Exemple 3} Passage d'un atlas \inserted{$\unsure{M|\mathcal{M}}$} à un sous-atlas \inserted{$M'$}, ou d'un magasin à un sous-magasin, dans l'esprit de restreindre le type de multistrates envisagées (par des propriétés du type connexité ou \unsure{inconnexité}, \uline{régularité}, être \deleted{\unsure{de}} des ``boules'' ou des ``simplexes'' dans un sens plus ou moins fort\unsure{''}) mais qui restent assez nombreuses pour ``\inserted{pouvoir} se \unsure{réduire} à $M'$'' -- notamment avec des \inserted{$M'$-}subdivisions \unsure{cofinales}.

Dans un tel cas, sûrement on demandera que pour le plongement $M'\longrightarrow M$, toutes les structures
\[
\leq,\quad\overset{\circ}{\ll},\quad\mathrel{|\circ|}\quad/\!/\quad\unsure{\gtrless|\bowtie},\quad\parallel
\]
soient conservées, à l'exception \unsure{tout} au plus de $\overset{\circ}{\ll}$ qui serait remplacé par $\underset{\mathrm{pol}}{\overset{\circ}{\ll}}$, et par conséquent $\gtrless$ par $\underset{\inserted{\mathrm{pol}}}{\gtrless}\deleted{\unsure{\mathrm{pol}}}$.

Donc on est ici dans le cas d'une notion d'hom \uline{de magasins}, qui est un plongement -- donc pas de difficultés du type de celles évoquées avec les \unsure{exemples|exemple} 1.
